% ============================================================
% Step4 教程 · Part II：工程实现详解
% 由 step4_tutorial.tex \input 引入（共用同一份导言区与主题色）
% ============================================================

\part{工程实现详解}

\begin{tipbox}{这一部分写什么}
仓库里真实做出来的 Step4 工程（RTL 共 1231 行：字库 684 + osd\_track 313 +
vision\_top 234）逐模块讲透，外加三次联调迭代的完整复盘。
对应目录：\code{03\_vision/step4\_osd\_integration/}。
\end{tipbox}

% ============================================================
\section{工程总览}

\subsection{文件清单}

\begin{center}
\begin{tabular}{llr l}
\toprule
\rowcolor{rowtint}\textbf{文件} & \textbf{类型} & \textbf{行数} & \textbf{职责} \\
\midrule
\code{font\_rom\_8x16.v} & 生成 & 684 & 40 字符 ASCII 子集点阵（脚本生成） \\
\code{osd\_track.v} & RTL & 313 & OSD 叠加层 v3（三行文本+双框+双缓冲） \\
\code{vision\_top.v} & RTL & 234 & 整机顶层 v2（检测+OSD+云台+声光） \\
\code{tb/tb\_osd\_track.v} & 测试 & 126 & OSD 整帧逐像素比对 \\
\code{tb/tb\_vision\_top.v} & 测试 & 99 & 整机 5 断言 \\
\code{tb/tb\_closed\_loop.v} & 测试 & 95 & 闭环在环收敛 \\
\code{sim\_out/*.png} & 证据 & --- & OSD 渲染对比图（金标准 vs RTL） \\
\bottomrule
\end{tabular}
\end{center}

\subsection{三时钟域与两条跨域河}

\begin{Verilog}
pix 域 (74.25MHz):  检测链 / OSD / 显示路径
pwm 域 (50MHz):     gimbal_ctrl / pwm_servo
跨域信号:
   frame_tick    pulse_sync (单 bit)
   23bit 坐标包  cdc_4phase
   135bit 目标包 cdc_4phase (同域备用, 接口兼容异域)
\end{Verilog}

\begin{keybox}{一条贯穿全项目的设计规则}
任何跨域信号只能走 pulse\_sync 或 cdc\_4phase 两座桥，禁止手搓两级同步器。
本步的 135bit 包就是 \code{cdc\_4phase \#.WIDTH(135)} 的一次参数化复用——
Step1 写好的桥，Step4 直接搬。
\end{keybox}

% ============================================================
\section{模块详解：font\_rom\_8x16（生成文件）}

字库 ROM 覆盖 OSD 用到的 40 个 ASCII 字符（大写字母、数字、空格、
\code{- + \% . :} 等）。生成链路：

\begin{Verilog}
cv2.putText(Hershey 字体) -> 8x16 阈值化点阵 -> case 表 RTL
地址 {char[7:0], row[3:0]} -> 该行 8bit (MSB = 最左像素)
\end{Verilog}

为什么选 8x16：16 像素高在 720p 上恰好是 45 行文字的量级，
评委在演示视频里看得清；8 像素宽让字符地址 = 像素 x 右移 3 位，
判决逻辑零乘法。

% ============================================================
\section{模块详解：osd\_track v3}

\subsection{v1 到 v3 的演进}

\begin{center}
\begin{tabular}{lp{8.6cm}}
\toprule
\rowcolor{rowtint}\textbf{版本} & \textbf{内容} \\
\midrule
v1 & 单行文字 + 准星 + 单框，68bit 包（IF-VISION v1.0.0） \\
v2 & 配合 Step5 拓展：包扩到 135bit，三行文本 + 跑马灯 + 运动框 \\
v3 & 消撕裂：影子寄存器组 + sof 拍整组换装（vsync 语义） \\
\bottomrule
\end{tabular}
\end{center}

\subsection{影子寄存器组：撕裂的结构性根治}

v3 的核心是\textbf{两组一模一样的寄存器}：

\begin{itemize}
  \item \textbf{影子组}（\code{s\_*}）：\code{pack\_valid} 拍写入，
    时刻可能在帧中间；
  \item \textbf{绘制组}（无前缀）：只在 \code{I\_valid \&\& I\_sof} 拍
    从影子整组拷贝，帧内再也不会变。
\end{itemize}

绘制判决只读绘制组——无论包何时到，一帧之内的坐标/文字/框
\textbf{是同一份快照}。已知残留：sof 拍本身（像素 (0,0)）仍用旧值绘制
一拍，但 OSD 起始在 x=16，无视觉影响（RTL 注释里明确记了这笔账）。

\subsection{文本构建：包到达拍的一次性算术}

\begin{Verilog}
txt1_s[0..2] = pack[133] ? "EDG" : "TRK"      // 算法名
txt1_s[3]    = 48 + preset                     // 阈值组号
txt1_s[5..8] = pack[134] ? "LOCK" : "LOST"    // 锁定态
txt1_s[11..14] = to_digits4(cx) 或 "----"      // 坐标四位
txt2: "D" lat " B" ±bb " C" cc " M" mmmm       // 遥测行
\end{Verilog}

延迟显示 \code{lat = disp\_cnt - fid}：包里带的是\textbf{采集时刻}的帧号，
显示端用当前帧号一减，就是这套链路的真实延迟帧数（进阶（3）的可视化）。

\subsection{跑马灯的算术}

\begin{Verilog}
mar_off = disp_cnt * 2                  // 每帧 2px
mar_pos = px - TXT_X0 + mar_off         // 相对位置 + 滚动偏移
mar_ci  = mar_pos[8:3] % 63             // 字符索引(63 字符取模)
on_mar  = 字形位[7 - mar_pos[2:0]]      // 字节内位选
\end{Verilog}

63 字符 x 8 像素 = 504 像素周期，2px/帧 即 252 帧（约 4.2 秒）滚一圈。

\subsection{准星与双框：带符号差 + 边缘四线}

\begin{Verilog}
dx_c = $signed(px) - $signed(cx);   // 带符号差, 一切判决的基础
ax = (|dx_c| <= 12) && (|dx_c| >= 3);   // 臂: 3..12 (中心 3px 留空)
on_cross = tv && ((ay && dx_c==0) || (ax && dy_c==0));
// 框: 四条边各 2px 厚
x_edge = (px==minx) || (px==minx+1) || (px==maxx-1) || (px==maxx);
on_box = tv && ((x_edge && in_by) || (y_edge && in_bx));
\end{Verilog}

\begin{warnbox}{带符号差是框类判决的唯一正确姿势}
如果用无符号减法算 \code{px - cx}，目标左半边的像素差全部回绕成
巨大正数，左半边框直接消失——和 Step3 闭环跑飞是同一类根因。
本项目已两次被"无符号减法"咬，写进了方法论。
\end{warnbox}

\subsection{输出 mux 与颜色}

优先级一个表达式写完（文字 $>$ 准星 $>$ 目标框 $>$ 运动框 $>$ 底图）：

\begin{Verilog}
O_data <= on_txt1|on_txt2 ? 24'hFFFFFF :   // 文字 白
          on_mar   ? 24'h00FFFF :          // 跑马灯 青
          on_cross ? 24'h00FF00 :          // 准星 绿
          on_box   ? 24'hFF0000 :          // 目标框 红
          on_mbox  ? 24'hFF8000 : I_data;  // 运动框 橙
\end{Verilog}

% ============================================================
\section{模块详解：vision\_top v2}

\subsection{算法切换：帧边界平稳换源}

进阶（1）要求多算法切换。工程做法：\textbf{边缘检测常开}（并行分支），
切换只动一个 mux 的选择端，且\textbf{在帧首拍锁存}：

\begin{Verilog}
if (I_tvalid && I_tuser)          // 只在帧首采样
    algo_q <= I_algo_edge;
// 之后一整帧 mux 固定选检测流或边缘流
\end{Verilog}

帧中间切会怎样？mux 前两条流水线延迟不同，切换拍前后像素来自
\textbf{两条链的不同进度}，画面错位——帧边界锁存让切换永远发生在
"两帧之间"，视觉无缝。

\subsection{135bit 包拼装与保留位}

\begin{Verilog}
wire [134:0] pack_vision = {t_valid, I_algo_edge, I_preset_idx,
    t_cx, t_cy, t_minx, t_maxx, t_miny, t_maxy, t_fid,
    mo_valid, mo_cnt[13:0], mo_minx, mo_maxx, mo_miny, mo_maxy,
    1'b0};    // <- bit0 保留位, 必须补 0
\end{Verilog}

RTL 注释里直接写着事故记录："缺位时整包上移一位，OSD 解析 valid 恒 0
（永远 LOST 事故根因，2026-09-20 修）"。\textbf{把教训写在案发现场}，
下一个改包的人第一眼就能看到。

\subsection{控制路径：缓存 + 帧拍消费}

23bit 控制包（valid+cx+cy）跨域后进 \code{ctrl\_hold} 缓存，
\code{ctrl\_vld\_hold} 标记有效；帧拍 \code{tick\_pwm} 到达时消费一次。
这样"包先到/帧拍先到"两种顺序都能正确工作，
代价是控制滞后 1 帧——设计行为，闭环增益公式里已经算过它。

\subsection{声光反馈（给嵌入式方向的接口）}

\code{led\_buzzer}：锁定亮绿灯、丢失红灯 1Hz 闪、\textbf{锁定沿}（从丢到锁）
蜂鸣 200ms。演示时评委不用盯着屏幕，听声音就知道系统锁上了。

% ============================================================
\section{测试平台设计}

\subsection{tb\_osd\_track：整帧逐像素}

激励：\code{osd\_coord.txt} 给出每帧的目标坐标/包围框（金标准写的），
TPG 底图 + OSD 叠加后整帧转储；期望：\code{osd\_exp*.hex} 是
Python \code{render\_osd()} 渲染的同帧画面。比对粒度：\textbf{像素级
24bit 相等}，4 帧共 3686400 个样本。

\subsection{tb\_vision\_top：五条整机断言}

\begin{center}
\begin{tabular}{ll}
\toprule
\rowcolor{rowtint}\textbf{断言} & \textbf{意义} \\
\midrule
直通输出与输入逐拍一致 & 主线零扰动（显示不受影响） \\
第 4 帧 OSD 与金标准一致 & 整链绘制的像素级正确 \\
duty\_x 在 [700,1500) & 云台 X 朝目标方向动 \\
duty\_y == 1500 & Y 轴死区正确（误差 20 = DZ\_Y 边界内） \\
lost\_cnt == 0 & 无误丢帧 \\
\bottomrule
\end{tabular}
\end{center}

\begin{warnbox}{迭代 3：采样过早的教训}
第一版在第 2 帧末读 duty，恒 1500 判 FAIL——其实云台正常，只是
"检测(1帧)+除法+跨域+消费"共约 2 帧滞后是设计行为。
修复：采样点移到第 4 帧末。\textbf{给断言选时机，要按数据通路的
真实延迟账来}，不能拍脑袋。
\end{warnbox}

\subsection{tb\_closed\_loop：相机模型驱动 TPG}

TPG 方块位置由 TB 相机模型实时计算：
\code{box\_x = 640 + (40 度 - theta) x 20px/度}——云台转多少，
方块就往回挪多少，\textbf{和真实相机的行为同构}。60 帧判据：
误差进入 ±25px 死区并保持 10 帧。实测 \code{settle\_frame=12}
（T1 控制节拍修正后；修正前为 36），
逐帧轨迹在 \code{sim/closed\_loop\_trace.txt}（答辩可直接画收敛曲线）。

% ============================================================
\section{验证结果与联调复盘}

\subsection{回归矩阵}

\begin{center}
\begin{tabular}{lll}
\toprule
\rowcolor{rowtint}\textbf{测试} & \textbf{比对粒度} & \textbf{结果} \\
\midrule
tb\_osd\_track & 整帧逐像素（4 帧 3686400 样本） & PASS \\
tb\_vision\_top & 5 断言 + 整帧 & PASS \\
tb\_closed\_loop & 在环收敛（settle\_frame） & PASS：12 帧 \\
\bottomrule
\end{tabular}
\end{center}

\subsection{"永远 LOST"事故复盘（本步骤最大缺陷）}

\begin{warnbox}{一次事故，三个裂缝}
现象：OSD 永远显示 LOST、云台 duty 恒 1500 不动、文字区坐标被污染。\\[3pt]
\textbf{根因链}：vision\_top 拼 69bit - 截成 67bit（valid 位丢）；
TB 侧 67bit 接 68bit 端口（MSB 悬空为 Z）；osd\_track 按 67bit 解包。
\textbf{三处三个宽度}，每一个单看都"能用"。\\[3pt]
\textbf{连带发现}：位偏移重排后，文字分支还读旧位 [66]（现在是 preset 位）
而不是 [67]（valid 位）——所以出现"显示 LOST 但坐标数字正常"
的灵异现象（数字用的是 [67]）。一字符修复。\\[3pt]
\textbf{修复动作}：统一宽度、位偏移重排、\textbf{冻结进 IF-VISION.md}
（接口唯一权威定义），此后任何位宽变更必须走接口变更记录。
\end{warnbox}

\begin{keybox}{为什么逐像素比对是抓这类 bug 的唯一手段}
包宽错误不会让仿真崩溃——valid 恒 0 的表现只是"画面对但没框"，
肉眼完全可能当成"场景里没有目标"。只有逐像素比对才把"应有框 vs 实际无框"
的每个像素差摆上台面。Step4 的三个真缺陷（包宽、旧位、字段溢出）
全部由它定位。
\end{keybox}

\subsection{Step4 真缺陷清单}

\begin{center}
\begin{tabular}{llp{5.6cm}l}
\toprule
\rowcolor{rowtint}\textbf{编号} & \textbf{现象} & \textbf{根因} & \textbf{类型} \\
\midrule
O1 & 永远 LOST & 包宽三处不一致（69/67/68） & 接口 \\
O2 & LOST 但坐标正常 & word 分支读旧位 [66] & RTL \\
O3 & 框画到 (72,18) & 坐标文件 \%02X 写 300 溢出 3 字符 & TB 脚本 \\
O4 & duty 恒 1500 误报 & TB 采样点早于通路真实延迟 & TB \\
\bottomrule
\end{tabular}
\end{center}

% ============================================================
\section{上板挂载说明（给上板组）}

\begin{center}
\begin{tabular}{ll}
\toprule
\rowcolor{rowtint}\textbf{项目} & \textbf{内容} \\
\midrule
挂载方式 A & 显示侧：官方 hdmi\_mixer 之前串入 osd\_track（vpix 域） \\
挂载方式 B & 检测侧：串在检测分支出口（本工程仿真即此形态） \\
vpix/DE 转换 & 显示侧若只有 DE 信号，按 IF-OSD.md 的规则生成 sof/eol \\
按键 & preset\_idx（阈值组）/ track\_en / algo\_edge / 亮度对比度 \\
观测 & ILA 挂 135bit 包 + O\_status；duty\_dbg 看云台 \\
\bottomrule
\end{tabular}
\end{center}

\begin{tipbox}{答辩演示建议}
画面里的三行 OSD 本身就是"评分表逐条自证"：L1 证明扩展（1）叠加与
进阶（2）坐标标注，D 字段证明进阶（3）延迟可视化，M 字段证明
进阶（4）运动检测，B/C 字段证明扩展（4）实时调节——
不需要额外口头解释。
\end{tipbox}

% ============================================================
\section{小结：Step4 的可复用产出}

\begin{itemize}
  \item \textbf{整机顶层 vision\_top}：检测/OSD/云台/声光四合一，
    上板即用；
  \item \textbf{一套 OSD 叠加层}：三行文本、双框、跑马灯、双缓冲，
    全部逐像素验证；
  \item \textbf{逐像素金标准比对方法}：此后任何绘制改动都有回归兜底；
  \item \textbf{接口治理范例}：135bit 包唯一权威定义 + 变更记录制度；
  \item \textbf{收敛证据链}：closed\_loop\_trace.txt 逐帧轨迹，
    答辩直接可画。
\end{itemize}

\begin{okbox}{下一步（Step5 预告）}
基础与进阶都通了，Step5 把评分表剩下的扩展项逐个补齐：
多算法切换、运动检测、亮度对比度、声光反馈——
每个都带位精确验证，并最终把回归扩到 16 个测试全绿。
\end{okbox}
